% భాగం: ప్రతిజ్ఞావాక్య తర్కం
% అధ్యాయం: ప్రతిజ్ఞావాక్య తర్కం
% విభాగం: నిర్దుష్టత

\documentclass[../../../include/open-logic-section]{subfiles}

\begin{document}

\olfileid{pl}{prp}{snd}

\olsection{ప్రతిజ్ఞావాక్య తర్కపు నిర్దుష్టత}

\begin{thm}[నిర్దుష్టత]
\ollabel{thm:soundness}
$\Gamma\Proves!A$ అయితే $\Gamma \Entails !A$.
\end{thm}

\begin{proof}
$\Gamma$ నుంచి $!A$కు ఉన్న వ్యుత్పత్తి పొడవుపై ఆగమనంతో నిరూపిస్తాం. $!A$ ఒక స్వీకృతమైతే $\Entails !A$. కాబట్టి $\Gamma \Entails!A$ కూడా సిద్ధిస్తుంది. అలాగే $!A \in \Gamma$ అయితే $\Gamma \Entails!A$. ఆగమన దశలో $!C$, $!C\lif!A$ నుంచి \emph{మోడస్ పోనెన్స్} ద్వారా $!A$ వచ్చిందనుకుందాం. $!C\lif !A$, $!C$లకు ఈ ఫలితం వర్తిస్తుందని ఆగమన పరికల్పనగా తీసుకుందాం. \olref[sem]{prop:semanticalfacts}లోని మూడో అంశం నుంచి కావలసిన ఫలితం వస్తుంది.
\end{proof}

\begin{cor}
$\Gamma$ సంతృప్తిపరచదగినదైతే, $\Gamma$ సుసంగతం. అందువల్ల ప్రతిజ్ఞావాక్య తర్కం సుసంగతం.
\end{cor}

\end{document}
